% 无向边与一对反向弧；同一三点连接（1—2—3）的无向形式与两种定向。
\begin{tikzpicture}[>=Stealth,
  v/.style={circle,draw=FigureStroke,fill=FigureBlueFill,line width=.8pt,minimum size=6.5mm,inner sep=0pt},
  e/.style={draw=FigureStroke,line width=.8pt},
  a/.style={draw=FigureStroke,line width=.8pt,->,shorten >=1pt},
  note/.style={font=\FigureNoteFont\fontsize{7.5}{9}\selectfont,text=FigureMuted,align=center}]
% 上排
\begin{scope}
  \node[v] (u) at (0,0) {$u$}; \node[v] (w) at (1.6,0) {$v$};
  \draw[e] (u)--(w);
  \node at (.8,-.85) {(a) 无向边$\{u,v\}$};
\end{scope}
\begin{scope}[xshift=3.4cm]
  \node[v] (u) at (0,0) {$u$}; \node[v] (w) at (1.6,0) {$v$};
  \draw[a] (u) to[bend left=25] (w);
  \draw[a] (w) to[bend left=25] (u);
  \node at (.8,-.85) {(b) 两条弧$(u,v)$与$(v,u)$};
\end{scope}
% 下排
\begin{scope}[yshift=-2.1cm]
  \node[v] (1) at (0,0) {$1$}; \node[v] (2) at (1.1,0) {$2$}; \node[v] (3) at (2.2,0) {$3$};
  \draw[e] (1)--(2)--(3);
  \node at (1.1,-.85) {(c) 无向：两两可达};
\end{scope}
\begin{scope}[yshift=-2.1cm,xshift=3.4cm]
  \node[v] (1) at (0,0) {$1$}; \node[v] (2) at (1.1,0) {$2$}; \node[v] (3) at (2.2,0) {$3$};
  \draw[a] (1)--(2); \draw[a] (2)--(3);
  \node at (1.1,-.85) {(d) $1$可到$3$，$3$不能回$1$};
\end{scope}
\begin{scope}[yshift=-2.1cm,xshift=6.8cm]
  \node[v] (1) at (0,0) {$1$}; \node[v] (2) at (1.1,0) {$2$}; \node[v] (3) at (2.2,0) {$3$};
  \draw[a] (1)--(2); \draw[a] (3)--(2);
  \node at (1.1,-.85) {(e) $1,3$都只能到$2$};
\end{scope}
\end{tikzpicture}
